\chapter{Реализация и экспериментальное исследование системы}\label{ch:ch3}

В главе описаны программная реализация спроектированной в
главе~\ref{ch:ch2} системы и результаты её экспериментального исследования на
реальной выборке электронной почты. Изложение ведётся по~этапам конвейера
обработки: для каждого этапа приводится ключевой фрагмент реализации и~получаемый
результат. \ph{черновая глава: числа и рисунки — плейсхолдеры}

\section{Среда реализации и организация кода}\label{sec:ch3/impl}

Система реализована на~языке Python (версии~3.11 и~выше). Прикладная логика
опирается на~набор внешних библиотек: \texttt{sentence-transformers} для~построения
векторных представлений моделью \texttt{bge-m3}~\cite{Chen2024BGEM3},
\texttt{hdbscan} для~плотностной кластеризации~\cite{McInnes2017HDBSCAN},
\texttt{optuna} для~подбора гиперпараметров~\cite{Akiba2019Optuna},
\texttt{scikit-learn} для~расчёта метрик и~вспомогательных алгоритмов~\cite{Pedregosa2011sklearn},
\texttt{torch} для~обучения и~применения нейросетевого
классификатора~\cite{Paszke2019PyTorch}, \texttt{psycopg} для~доступа к~СУБД
PostgreSQL и~\texttt{argilla} для~организации интерфейса ручной
разметки~\cite{Vila2023Argilla}.

Исходный код организован в~пакете \texttt{diplom\_ai} и~разбит на~слои в~соответствии
с~архитектурой, описанной в~главе~\ref{ch:ch2}; назначение пакетов сведено
в~таблицу~\ref{tab:ch3/modules}. Каждый слой отвечает за~отдельный этап конвейера
обработки и~взаимодействует со~смежными слоями через явно заданные интерфейсы, что
упрощает модульное тестирование и~замену реализаций.

Качество кода поддерживается набором инструментов: \texttt{ruff} (статический
анализ и~форматирование), \texttt{mypy} (проверка типов) и~\texttt{pytest}
(модульные тесты). Инфраструктурные зависимости (PostgreSQL, сервис разметки
Argilla) разворачиваются единообразно через Docker Compose, что обеспечивает
воспроизводимость окружения у~разных исполнителей.

\begin{table}[htbp]
    \centering
    \caption{Пакеты приложения и их назначение}\label{tab:ch3/modules}
    \begin{tabular}{@{}ll@{}}
        \toprule
        Пакет & Назначение \\
        \midrule
        \texttt{email}         & синхронизация и нормализация писем \\
        \texttt{storage}       & доступ к PostgreSQL, схема, репозитории \\
        \texttt{queueing}      & очередь задач конвейера \\
        \texttt{orchestration} & диспетчеризация этапов \\
        \texttt{clustering}    & эмбеддинги, кластеризация, подбор параметров \\
        \texttt{annotation}    & экспорт/импорт разметки (Argilla) \\
        \texttt{training}      & обучение классификатора \\
        \texttt{inference}     & применение модели \\
        \texttt{evaluation}    & метрики качества \\
        \bottomrule
    \end{tabular}
\end{table}

Этапы конвейера выполняются не~напрямую, а~через очередь задач: оркестратор
ставит в~очередь единицы работы (например, нормализацию каждого письма) и~затем
дренирует очередь, последовательно исполняя задачи с~обработкой повторных попыток
при~сбоях. Такой подход обеспечивает идемпотентность и~устойчивость к~частичным
отказам. Листинг~\ref{lst:ch3/drain} показывает реализацию этого механизма
на~примере этапа нормализации.

\lstinputlisting[language={Python},linerange={112-156},%
    caption={Постановка задач этапа в очередь и её дренаж},label={lst:ch3/drain}]%
    {listings/email_pipeline.py}

\section{Данные, экспериментальная установка и протокол оценки}\label{sec:ch3/setup}

Экспериментальное исследование выполнено на~обезличенной выборке писем реального
почтового ящика. Из~соображений конфиденциальности приводятся только агрегированные
характеристики корпуса без~раскрытия содержания отдельных сообщений. После
синхронизации и~нормализации выборка содержит \ph{N} писем; на~этапе оценки качества
сообщения распределяются по~категориям \texttt{ok}, \texttt{empty} и~\texttt{too\_short},
а~также по~языку (русский / английский). Сводный состав выборки приведён
в~таблице~\ref{tab:ch3/corpus}.

\begin{table}[htbp]
    \centering
    \caption{Состав исследуемой выборки писем}\label{tab:ch3/corpus}
    \begin{tabular}{@{}lr@{}}
        \toprule
        Характеристика & Значение \\
        \midrule
        Всего писем (после синхронизации) & \ph{N} \\
        Качество \texttt{ok} / \texttt{empty} / \texttt{too\_short} & \ph{..\,/\,..\,/\,..} \\
        Доля русскоязычных / англоязычных & \ph{0.NN} / \ph{0.NN} \\
        Период & \ph{ММ.ГГГГ--ММ.ГГГГ} \\
        \bottomrule
    \end{tabular}
\end{table}

Конфигурация эксперимента фиксирована для~воспроизводимости. Векторные представления
строятся моделью \texttt{BAAI/bge-m3} (версия~\ph{..}); все этапы, использующие
псевдослучайность (подбор гиперпараметров, разбиение на~обучающую и~тестовую части,
инициализация классификатора), запускаются с~фиксированным начальным значением
генератора \ph{seed}.

Качество результатов оценивается метриками, введёнными в~главе~\ref{ch:problem}. Для~оценки
кластеризации используются внутренние меры из~раздела~\ref{subsec:ch1/cluster/metrics}:
индекс достоверности плотностной кластеризации (DBCV), силуэтный коэффициент и~индекс
Дэвиса--Боулдина. Для~оценки классификатора применяются метрики
раздела~\ref{subsec:ch1/task/metrics}: доля верных ответов (accuracy), макроусреднённая
и~взвешенная меры \( F_1 \), а~также покрытие и~доля писем, отнесённых к~служебной
категории \texttt{other}.

Классификатор оценивается на~золотом тестовом множестве, сформированном из~человеко-проверенных
меток. Чтобы исключить утечку информации, возникающую при~распространении метки
с~представителя на~весь кластер, разбиение на~обучающую и~тестовую части выполняется
\emph{по~кластерам}: письма одного кластера целиком попадают либо в~обучение, либо
в~тест, но~не~в~оба множества одновременно. Разбиение детерминировано при~фиксированном
\ph{seed}; для~каждого класса дополнительно фиксируется число тестовых примеров
(support), что позволяет корректно интерпретировать поклассовые показатели.

\section{Синхронизация и нормализация писем}\label{sec:ch3/norm}

Письма проходят путь \texttt{IMAP}~\(\to\)~сырое MIME-сообщение~\(\to\)~запись
\texttt{RawEmail}~\(\to\)~нормализованная запись \texttt{NormalizedEmail}. Базовая
токенизация текста реализована в~модуле \texttt{data/preprocessing.py} функцией
\texttt{normalize\_text}: ссылки, адреса электронной почты, даты, время и~числа
заменяются на~устойчивые токены \texttt{LINK}, \texttt{EMAIL}, \texttt{DATE},
\texttt{TIME}, \texttt{NUM}, после чего схлопываются пробельные символы
(листинг~\ref{lst:ch3/tokens}). Такая замена снижает разреженность словаря и~убирает
несущественную для~тематической классификации вариативность.

\lstinputlisting[language={Python},linerange={5-21},%
    caption={Замена ссылок, адресов, дат и чисел токенами},label={lst:ch3/tokens}]%
    {listings/preprocessing.py}

Модуль \texttt{email/normalization.py} (функция \texttt{normalize\_raw\_email})
не~дублирует токенизацию, а~\emph{делегирует} её \texttt{normalize\_text},
дополнительно выполняя извлечение основного текста из~MIME-структуры, срез цитат,
подписей и~футеров, определение языка и~оценку качества полученного текста
(листинг~\ref{lst:ch3/normalize}).

\lstinputlisting[language={Python},linerange={41-55},%
    caption={Нормализация письма: делегирование токенизации и срез цитат},%
    label={lst:ch3/normalize}]{listings/normalization.py}

В результате каждое письмо получает признак языка и~категорию качества. Сообщения
категорий \texttt{empty} и~\texttt{too\_short} исключаются из~дальнейшей обработки как
не~несущие тематической информации. Распределение по~качеству и~языку приведено
в~таблице~\ref{tab:ch3/norm-quality}; типичный результат нормализации показан
фрагментом «до~/~после»: \ph{обезличенный пример «до» $\to$ «после»}.

\begin{table}[htbp]
    \centering
    \caption{Результат нормализации: качество и язык}\label{tab:ch3/norm-quality}
    \begin{tabular}{@{}lrr@{}}
        \toprule
        Категория & Русские & Английские \\
        \midrule
        \texttt{ok}        & \ph{..} & \ph{..} \\
        \texttt{empty}     & \ph{..} & \ph{..} \\
        \texttt{too\_short}& \ph{..} & \ph{..} \\
        \bottomrule
    \end{tabular}
\end{table}

\section{Построение векторных представлений}\label{sec:ch3/embed}

Нормализованные письма, прошедшие фильтр качества, преобразуются в~плотные
векторные представления моделью \texttt{bge-m3}~\cite{Chen2024BGEM3}. Кодирование
выполняется пакетами (батчами) фиксированного размера, что позволяет эффективно
использовать вычислительный ресурс; полученные векторы накапливаются и~объединяются
в~единую матрицу (листинг~\ref{lst:ch3/embeddings}).

\lstinputlisting[language={Python},linerange={149-167},%
    caption={Батчевое построение эмбеддингов и запись в кеш \texttt{.npz}},%
    label={lst:ch3/embeddings}]{listings/embeddings.py}

Матрица векторов сохраняется в~файл формата \texttt{.npz} атомарно: запись ведётся
во~временный файл с~последующим переименованием (\texttt{os.replace}), что исключает
повреждение кеша при~аварийном завершении. Для~каждого письма в~таблицу
\texttt{email\_embeddings} заносится индексная запись со~ссылкой на~файл, номером
строки в~матрице, а~также именем и~версией модели; версионирование позволяет
пересчитывать представления при~смене модели, не~затирая прежние. Повторный запуск
этапа пропускает письма, для~которых представление уже построено.

После фильтрации по~языку и~качеству в~кеш записано \ph{N} векторов размерности
\ph{d}; объём кеша составляет \ph{..}~МБ, суммарное время построения~--- \ph{..}.

\section{Кластеризация и подбор параметров}\label{sec:ch3/cluster}

Тематическая структура корпуса выявляется плотностной кластеризацией
HDBSCAN~\cite{Campello2013HDBSCAN,McInnes2017HDBSCAN}, применяемой к~векторным
представлениям писем. В~отличие от~партиционных методов, HDBSCAN не~требует
заранее заданного числа кластеров и~явно выделяет шумовые точки, что соответствует
природе почтового потока, где значительная доля писем не~образует плотных тематических
групп.

Гиперпараметры \texttt{min\_cluster\_size} и~\texttt{min\_samples} подбираются
байесовской оптимизацией Optuna~\cite{Akiba2019Optuna} (сэмплер TPE). Целевая функция
(листинг~\ref{lst:ch3/tuner}) максимизирует индекс DBCV~\cite{Moulavi2014DBCV}
с~добавками за~устойчивость кластеров и~уверенность принадлежности и~со~штрафом
за~избыточную долю шума; конфигурации, нарушающие ограничения допустимости (число
кластеров вне~заданного диапазона, доля шума выше порога, слишком мелкие кластеры),
отбраковываются возвратом штрафного значения.

\lstinputlisting[language={Python},linerange={36-97},%
    caption={Целевая функция подбора параметров HDBSCAN (DBCV с добавками)},%
    label={lst:ch3/tuner}]{listings/tuner.py}

Качество кластеризации сопоставлено с~двумя опорными конфигурациями: HDBSCAN
с~параметрами по~умолчанию и~партиционным методом KMeans~\cite{MacQueen1967KMeans,Arthur2007KMeansPP}.
Число кластеров KMeans задаётся равным числу кластеров, найденному настроенным
HDBSCAN. Силуэтный коэффициент~\cite{Rousseeuw1987Silhouette} и~индекс
Дэвиса--Боулдина~\cite{Davies1979DBI} вычисляются на~не-шумовых точках обоих
алгоритмов, индекс DBCV~--- только для~HDBSCAN (для~KMeans он неприменим, так как метод
не~выделяет шум), а~доля покрытия (единица минус доля шума) приводится как контекст.

\begin{figure}[htbp]
    \centering
    \phbox{рис.: история оптимизации Optuna --- заменить на \texttt{optuna\_history.pdf}}
    \caption{История оптимизации DBCV (Optuna, TPE)}\label{fig:ch3/optuna}
\end{figure}

\begin{table}[htbp]
    \centering
    \caption{Лучшие гиперпараметры HDBSCAN и важность параметров}\label{tab:ch3/optuna-best}
    \begin{tabular}{@{}lr@{}}
        \toprule
        Параметр & Значение / важность \\
        \midrule
        \texttt{min\_cluster\_size} & \ph{..} \\
        \texttt{min\_samples}       & \ph{..} \\
        Значение целевой функции     & \ph{0.NN} \\
        Важность \texttt{min\_cluster\_size} / \texttt{min\_samples} & \ph{0.NN} / \ph{0.NN} \\
        \bottomrule
    \end{tabular}
\end{table}

\begin{table}[htbp]
    \centering
    \caption{Сравнение алгоритмов кластеризации (внутренние метрики)}\label{tab:ch3/cluster-compare}
    \begin{tabular}{@{}lrrr@{}}
        \toprule
        Метрика & HDBSCAN+тюнер & HDBSCAN (дефолт) & KMeans \\
        \midrule
        Число кластеров            & \ph{..} & \ph{..} & \ph{..} \\
        Силуэт (без шума)          & \ph{0.NN} & \ph{0.NN} & \ph{0.NN} \\
        Индекс Дэвиса--Боулдина    & \ph{0.NN} & \ph{0.NN} & \ph{0.NN} \\
        DBCV                       & \ph{0.NN} & \ph{0.NN} & --- \\
        Доля покрытия (1 -- шум)   & \ph{0.NN} & \ph{0.NN} & 1.00 \\
        \bottomrule
    \end{tabular}
\end{table}

\begin{table}[htbp]
    \centering
    \caption{Метрики итогового запуска кластеризации}\label{tab:ch3/cluster-metrics}
    \begin{tabular}{@{}lr@{}}
        \toprule
        Показатель & Значение \\
        \midrule
        Число кластеров          & \ph{K} \\
        Доля шума                & \ph{0.NN} \\
        DBCV                     & \ph{0.NN} \\
        Силуэт                   & \ph{0.NN} \\
        Индекс Дэвиса--Боулдина  & \ph{0.NN} \\
        \bottomrule
    \end{tabular}
\end{table}

\begin{figure}[htbp]
    \centering
    \phbox{рис.: 2D-проекция кластеров (UMAP) --- заменить на \texttt{cluster\_projection.pdf}}
    \caption{Двумерная проекция кластеров (UMAP)}\label{fig:ch3/projection}
\end{figure}

Для~содержательной интерпретации найденных групп для~каждого кластера выделены
ключевые термины по~мере TF-IDF, а~их взаимное расположение визуализировано
двумерной проекцией методом UMAP~\cite{McInnes2018UMAP} (рис.~\ref{fig:ch3/projection});
сводка приведена в~таблице~\ref{tab:ch3/topterms}.

\begin{table}[htbp]
    \centering
    \caption{Ключевые термины кластеров (TF-IDF)}\label{tab:ch3/topterms}
    \begin{tabular}{@{}llr@{}}
        \toprule
        Кластер & Топ-термины & Размер \\
        \midrule
        1 & \ph{термины} & \ph{..} \\
        2 & \ph{термины} & \ph{..} \\
        \multicolumn{3}{@{}l}{\ph{... остальные кластеры}} \\
        \bottomrule
    \end{tabular}
\end{table}

\section{Ручная проверка и формирование классов}\label{sec:ch3/annot}

\section{Классификатор: обучение и оценка}\label{sec:ch3/clf}

\section{Применение классификатора и порог уверенности}\label{sec:ch3/infer}

\section{Демонстрация цикла переобучения}\label{sec:ch3/loop}

\section{Выводы по главе}\label{sec:ch3/concl}

\clearpage
